\chapter{Sprint 4 : Planification visuelle et notifications}

\section*{Introduction}
Le quatrième sprint apporte les outils de visualisation décisionnelle : diagramme de
Gantt interactif avec calcul automatique du chemin critique (algorithme CPM),
calendrier mensuel des tâches, système de notifications temps réel via WebSocket
(Laravel Reverb), et tableaux de bord différenciés selon le profil de l'utilisateur.

% ============================================================
\section{Sprint Backlog}

\begin{table}[H]
  \centering
  \caption{Sprint Backlog — Sprint 4}
  \label{tab:backlog_s4}
  \begin{tabularx}{\textwidth}{C{1.2cm}L{5.5cm}L{2cm}C{2cm}C{1.5cm}}
    \toprule
    \thead{ID} & \thead{User Story} & \thead{Acteur} & \thead{Priorité} & \thead{Points} \\
    \midrule
    \rowcolor{lightgray}
    US-17 & Visualiser le Gantt avec chemin critique     & Member & Must   & 13 \\
    US-18 & Consulter le calendrier mensuel              & Member & Should &  5 \\
    \rowcolor{lightgray}
    US-19 & Recevoir des notifications temps réel        & Tous   & Must   &  8 \\
    US-20 & Configurer ses préférences de notification   & Tous   & Should &  3 \\
    \rowcolor{lightgray}
    US-21 & Accéder à un dashboard personnel             & Member & Should &  8 \\
    US-22 & Accéder au dashboard administrateur          & Admin  & Should &  5 \\
    \rowcolor{lightgray}
    US-27 & Consulter le rapport d'avancement (client)   & Client & Should &  5 \\
    US-28 & Utiliser la recherche globale                & Tous   & Should &  5 \\
    \midrule
    \multicolumn{4}{r}{\textbf{Total :}} & \textbf{52} \\
    \bottomrule
  \end{tabularx}
\end{table}

% ============================================================
\section{Conception}

\subsection{Diagramme de cas d'utilisation}

\begin{figure}[H]
  \centering
  \IfFileExists{figures/sprint4_use_case.png}{%
    \includegraphics[width=0.88\textwidth]{figures/sprint4_use_case.png}%
  }{%
    \fbox{\parbox{0.88\textwidth}{\centering\vspace{1.8cm}
      \textbf{[IMAGE A INSÉRER]}\\[0.4em]
      \texttt{figures/sprint4\_use\_case.png}\\[0.3em]
      \textit{Cas d'utilisation : Admin, Member, Client — Gantt, calendrier, notifications, dashboards}
    \vspace{1.8cm}}}%
  }
  \caption{Diagramme de cas d'utilisation — Sprint 4}
  \label{fig:uc_s4}
\end{figure}

\subsection{Diagramme de séquence — Détection de retard et notification}

\begin{figure}[H]
  \centering
  \IfFileExists{figures/sprint4_sequence_notif.png}{%
    \includegraphics[width=0.90\textwidth]{figures/sprint4_sequence_notif.png}%
  }{%
    \fbox{\parbox{0.90\textwidth}{\centering\vspace{1.8cm}
      \textbf{[IMAGE A INSÉRER]}\\[0.4em]
      \texttt{figures/sprint4\_sequence\_notif.png}\\[0.3em]
      \textit{Séquence : Scheduler horaire → détection retard → NotificationService → Reverb WebSocket}
    \vspace{1.8cm}}}%
  }
  \caption{Diagramme de séquence — Notification de retard (scheduler horaire)}
  \label{fig:seq_notif}
\end{figure}

Le scheduler Laravel s'exécute toutes les heures : il sélectionne les tâches dont la
\texttt{due\_date} est dépassée et dont le \texttt{workflow\_status} n'est pas
\texttt{done}. Pour chaque tâche concernée, \texttt{NotificationService::notifyDelay()}
persiste une notification en base et diffuse l'événement \texttt{TaskDelayed}
sur le canal privé \texttt{private-user.\{id\}}. Côté React, le badge de la cloche
s'incrémente et un toast apparaît.

\subsection{Diagramme de classes}

\begin{figure}[H]
  \centering
  \IfFileExists{figures/sprint4_class.png}{%
    \includegraphics[width=0.85\textwidth]{figures/sprint4_class.png}%
  }{%
    \fbox{\parbox{0.85\textwidth}{\centering\vspace{1.8cm}
      \textbf{[IMAGE A INSÉRER]}\\[0.4em]
      \texttt{figures/sprint4\_class.png}\\[0.3em]
      \textit{Classes : Notification, ActivityLog, CriticalPathService, AdminDashboardService}
    \vspace{1.8cm}}}%
  }
  \caption{Diagramme de classes — Sprint 4}
  \label{fig:class_s4}
\end{figure}

% ============================================================
\section{Réalisation}

\subsection{Diagramme de Gantt}

\begin{figure}[H]
  \centering
  \IfFileExists{figures/sprint4_19_gantt.png}{%
    \includegraphics[width=0.92\textwidth]{figures/sprint4_19_gantt.png}%
  }{%
    \fbox{\parbox{0.92\textwidth}{\centering\vspace{2cm}
      \textbf{[IMAGE A INSÉRER]}\\[0.4em]
      \texttt{figures/sprint4\_19\_gantt.png}\\[0.3em]
      \textit{Gantt : 12 projets janv.--déc. 2026, repère AUJOURD'HUI, couleurs par priorité,
      chemin critique en rouge dans la vue projet}
    \vspace{2cm}}}%
  }
  \caption{Diagramme de Gantt interactif — Vue globale multi-projets}
  \label{fig:ui_gantt}
\end{figure}

Le Gantt est rendu entièrement en SVG côté client. Chaque barre représente une tâche ;
sa couleur encode la priorité. Le repère vertical « AUJOURD'HUI » est calculé
dynamiquement. Dans la vue projet, les tâches du chemin critique (float = 0) sont
affichées en rouge. Un survol de barre affiche une infobulle avec titre, dates et
progression.

\subsection{Calendrier mensuel}

\begin{figure}[H]
  \centering
  \IfFileExists{figures/sprint4_20_calendar.png}{%
    \includegraphics[width=0.88\textwidth]{figures/sprint4_20_calendar.png}%
  }{%
    \fbox{\parbox{0.88\textwidth}{\centering\vspace{2cm}
      \textbf{[IMAGE A INSÉRER]}\\[0.4em]
      \texttt{figures/sprint4\_20\_calendar.png}\\[0.3em]
      \textit{Calendrier Mai 2026, filtres projet/utilisateur, événements colorés par priorité}
    \vspace{2cm}}}%
  }
  \caption{Calendrier mensuel — Mai 2026}
  \label{fig:ui_calendar}
\end{figure}

\subsection{Centre de notifications}

\begin{figure}[H]
  \centering
  \IfFileExists{figures/sprint4_21_notifications.png}{%
    \includegraphics[width=0.80\textwidth]{figures/sprint4_21_notifications.png}%
  }{%
    \fbox{\parbox{0.80\textwidth}{\centering\vspace{2cm}
      \textbf{[IMAGE A INSÉRER]}\\[0.4em]
      \texttt{figures/sprint4\_21\_notifications.png}\\[0.3em]
      \textit{Centre de notifications : onglets Tout / Invitations / Retards, badge compteur}
    \vspace{2cm}}}%
  }
  \caption{Centre de notifications — Onglets et filtres}
  \label{fig:ui_notifications}
\end{figure}

\subsection{Code clé — \texttt{CriticalPathService} et scheduler}

\begin{lstlisting}[style=phpstyle,caption={CriticalPathService::calculate() — Algorithme CPM},label={lst:cpm}]
// app/Services/CriticalPathService.php
public function calculate(Project $project): array
{
    $tasks = $project->tasks()->with('dependencies')->get();

    // --- Forward Pass ---
    foreach ($tasks->sortBy('topological_order') as $task) {
        $es = $task->dependencies
            ->map(fn($d) => $d->pivotTask->early_finish)->max() ?? 0;
        $task->early_start  = $es;
        $task->early_finish = $es + ($task->estimated_hours ?? 0);
    }

    $projectEnd = $tasks->max('early_finish');

    // --- Backward Pass ---
    foreach ($tasks->sortByDesc('topological_order') as $task) {
        $lf = $task->dependents
            ->map(fn($d) => $d->late_start)->min() ?? $projectEnd;
        $task->late_finish = $lf;
        $task->late_start  = $lf - ($task->estimated_hours ?? 0);
        $task->float       = $task->late_start - $task->early_start;
        $task->is_critical = ($task->float == 0);
    }

    return $tasks->filter->is_critical->values()->toArray();
}
\end{lstlisting}

\begin{lstlisting}[style=phpstyle,caption={Scheduler horaire — Détection des retards},label={lst:scheduler}]
// routes/console.php
Schedule::call(function () {
    $delayed = Task::whereDate('due_date', '<', now())
        ->whereNotIn('workflow_status', ['done'])
        ->where('updated_at', '<', now()->subHour())
        ->with('assignees')
        ->get();

    foreach ($delayed as $task) {
        app(NotificationService::class)
            ->notifyDelay($task, $task->assignees);
    }
})->hourly()->name('detect-delays')->withoutOverlapping();
\end{lstlisting}

% ============================================================
\section{Tests unitaires}

\begin{lstlisting}[style=phpstyle,caption={Tests unitaires — CPM et notifications},label={lst:tests_s4}]
it('identifies critical path tasks (float = 0)', function () {
    $project  = Project::factory()->withLinkedTasks(3)->create();
    $critical = app(CriticalPathService::class)->calculate($project);

    expect($critical)->not->toBeEmpty()
        ->and(collect($critical)->every(fn($t) => $t['float'] == 0))
        ->toBeTrue();
});

it('broadcasts delay notification', function () {
    Event::fake([TaskDelayed::class]);
    $task = Task::factory()->overdue()->create();

    app(NotificationService::class)->notifyDelay($task, $task->assignees);

    Event::assertDispatched(TaskDelayed::class);
});

it('marks notification as read', function () {
    $notif = Notification::factory()->unread()->create();
    $this->actingAs($notif->user)
         ->patchJson("/api/notifications/{$notif->id}/read")
         ->assertOk();
    expect($notif->fresh()->read_at)->not->toBeNull();
});
\end{lstlisting}

\section*{Conclusion}
Le Gantt avec chemin critique a été le composant le plus complexe du projet : la
sérialisation topologique du graphe de dépendances et le calcul CPM ont nécessité
un soin particulier pour gérer les cas limites (tâches sans dépendances, graphes
déconnectés). Laravel Reverb s'est montré très fiable, avec une latence médiane
inférieure à 20\,ms pour la diffusion des événements. Le système de dashboard
multi-rôles offre une expérience différenciée appréciée de tous les profils testés.
